\section{Related Work}

\subsection{LLM-based Access Control for IoT}

Recent work has explored large language models for IoT access control, primarily focusing on policy generation from natural language specifications.

\textbf{LACE}~\cite{lace2025} introduces a Language-based Access Control Engine that translates natural language policy descriptions into structured rules. LACE combines LLM-based policy generation with formal verification using Open Policy Agent (OPA) and achieves 88\% runtime decision accuracy with a 0.79 F1-score. However, LACE focuses on \textit{offline policy authoring}—users describe policies in natural language which are then compiled into OPA rules. In contrast, our system makes \textit{real-time runtime decisions} on individual access requests using live LLM inference.

\textbf{CollabIoT}~\cite{collabiot2025} employs LLM-driven capability-based access control for transient IoT collaboration. Their LLM pipeline generates validated access control tokens with 100\% policy correctness and achieves sub-millisecond enforcement overhead (0.3 ms). CollabIoT excels at \textit{policy generation} for device auto-configuration (150 ms setup time) but does not use LLMs for runtime access decisions. Our hybrid approach complements these systems by providing contextual runtime analysis when deterministic rules are insufficient.

\textbf{Key Distinction}: Both LACE and CollabIoT leverage LLMs for \textit{policy creation}—translating human intent into machine-enforceable rules. Our system addresses a different challenge: making \textit{real-time access decisions} on alerts using multi-agent LLM reasoning when deterministic validation cannot resolve semantic ambiguity. Table~\ref{tab:related_comparison} summarizes these differences.

\begin{table}[h]
    \centering
    \caption{Comparison with LLM-based IoT Access Control Systems}
    \label{tab:related_comparison}
    \begin{tabular}{lccc}
        \toprule
        \textbf{System} & \textbf{LLM Role} & \textbf{Decision Time} & \textbf{Accuracy} \\
        \midrule
        LACE~\cite{lace2025} & Policy generation & Offline & 88\% (runtime) \\
        CollabIoT~\cite{collabiot2025} & Token generation & Offline (150ms) & 100\% (policy) \\
        \textbf{Ours (Hybrid)} & Runtime decisions & Real-time (97ms) & \textbf{94.2\%} \\
        \bottomrule
    \end{tabular}
\end{table}

\subsection{Traditional Access Control for IoT}

Traditional IoT access control relies on static rule-based approaches such as RBAC (Role-Based Access Control)~\cite{rbac2004}, ABAC (Attribute-Based Access Control)~\cite{abac2012}, and capability-based systems~\cite{capability2016}. These methods excel at deterministic validation (user→device permissions, time-based policies) but struggle with semantic threats and contextual anomalies.

\textbf{DLBAC}~\cite{dlbac2022} proposes deep-learning-based access control using neural networks to learn access patterns from historical data. While improving over static RBAC in dynamic scenarios, DLBAC lacks explainability—a critical requirement for M0801 compliance where security decisions must be auditable.

\textbf{RASA}~\cite{rasa2021} introduces risk-aware adaptive access control using Bayesian networks to model uncertainty. However, RASA's risk scores are not human-interpretable, making it difficult for security operators to validate decisions.

Our hybrid approach preserves the reliability of deterministic validation for security-critical checks (M0801 user authentication) while augmenting it with LLM-based semantic reasoning that provides natural language explanations.

\subsection{Multi-Agent LLM Systems}

Recent work on multi-agent LLM systems~\cite{multiagent2023} demonstrates that specialized agents collaborating on a task outperform single general-purpose models. Our Policy Agent and Context Agent follow this pattern: the Context Agent extracts behavioral anomalies from alert history, while the Policy Agent evaluates access against defined policies. This separation of concerns improves both accuracy and explainability compared to monolithic prompts.
